% !TEX root = ../main.tex
\begin{ouverture}[gravitation]{16}{Mouvement dans un champ de gravitation}{Partie 2 · Mouvement et interactions}
{Satellites et planètes\enspace·\enspace Mouvement circulaire uniforme\enspace·\enspace Lois de Kepler\enspace·\enspace Satellite géostationnaire}
\objectif{Déterminer les vecteurs vitesse et accélération d'un satellite en orbite circulaire.}
\objectif{Établir et exploiter la troisième loi de Kepler dans le cas d'une orbite circulaire.}
\objectif{Énoncer et exploiter les trois lois de Kepler.}
\objectif{Caractériser l'orbite d'un satellite géostationnaire.}
\objectif{Tester les deuxième et troisième lois de Kepler en Python à partir de données astronomiques.}
\prerequis{Repère de Frenet, mouvement circulaire uniforme (chapitre~13).}
\prerequis{Deuxième loi de Newton, référentiels galiléens (chapitre~14).}
\prerequis{Loi de la gravitation universelle, champ de gravitation (classe de première).}
\end{ouverture}
\essentiel
\essSection{1}{Force et champ de gravitation}
\begin{loi}[Gravitation universelle]
Deux corps à répartition sphérique de masse, de masses $m$ et $M$, dont les centres sont distants de $r$, exercent l'un
sur l'autre des forces attractives de même valeur
\[F=G\,\frac{m\,M}{r^2},\qquad G=\qty{6,67e-11}{N.m^2.kg^{-2}}.\]
Le champ de gravitation créé par l'astre de masse $M$ à la distance $r$ de son centre est dirigé vers ce centre et vaut
$\mathcal{G}=G\,\dfrac{M}{r^2}$ : un corps de masse $m$ y subit la force $\vec{F}=m\,\vec{\mathcal{G}}$.
\end{loi}
\begin{exemple}
À la surface de la Terre ($M_\mathrm{T}=\qty{5,97e24}{kg}$, $R_\mathrm{T}=\qty{6371}{km}$),
$\mathcal{G}=\dfrac{\num{6.67e-11}\times\num{5.97e24}}{(\num{6.371e6})^2}=\qty{9,81}{N.kg^{-1}}$ : on retrouve l'intensité de la
pesanteur. À 420~km d'altitude, il vaut encore \qty{8,6}{N.kg^{-1}}.
\end{exemple}
\essSection{2}{Mouvement circulaire d'un satellite}
\begin{propriete}[Un mouvement circulaire est uniforme]
Un satellite de masse $m$, soumis à la seule attraction d'un astre de masse $M$, décrit une orbite circulaire de rayon
$r$. Dans le référentiel lié au centre de l'astre (géocentrique pour la Terre, héliocentrique pour le Soleil), supposé
galiléen, la deuxième loi de Newton s'écrit, dans le repère de Frenet :
\[m\,\vec{a}=G\,\frac{m\,M}{r^2}\,\un\qquad\text{soit}\qquad a_\mathrm{t}=\ddt{v}=0\quad\text{et}\quad\frac{v^2}{r}=\frac{G\,M}{r^2}.\]
La vitesse est constante : le mouvement est \textbf{circulaire uniforme}, et
\[v=\sqrt{\frac{G\,M}{r}},\qquad T=\frac{2\pi r}{v}=2\pi\sqrt{\frac{r^3}{G\,M}}.\]
Ces grandeurs ne dépendent pas de la masse du satellite.
\end{propriete}
\schema{16-orbite}
\begin{exemple}
Pour la Station spatiale internationale, $r=6371+420=\qty{6791}{km}$ :
$v=\sqrt{\dfrac{\num{3.98e14}}{\num{6.791e6}}}=\qty{7,66e3}{m.s^{-1}}$ et $T=\dfrac{2\pi r}{v}=\qty{5,57e3}{s}$, soit 92,9~minutes.
\end{exemple}
\essSection{3}{Les lois de Kepler}
\begin{loi}[Lois de Kepler]
Dans le référentiel héliocentrique :
\begin{itemize}[leftmargin=1.4em,itemsep=0.1em]
\item \textbf{première loi} : l'orbite d'une planète est une ellipse dont le Soleil occupe l'un des foyers ;
\item \textbf{deuxième loi} (loi des aires) : le segment qui relie le Soleil à la planète balaie des aires égales pendant
des durées égales ; la planète va donc plus vite près du Soleil ;
\item \textbf{troisième loi} : le rapport $\dfrac{T^2}{a^3}$ est le même pour toutes les planètes, $a$ étant le demi-grand axe de
l'orbite.
\end{itemize}
Ces lois s'appliquent à tout corps en orbite autour d'un astre : satellites de la Terre, lunes de Jupiter\,\dots
\end{loi}
\schema{16-kepler}
\begin{savoirfaire}[Établir la troisième loi pour une orbite circulaire]
Élever au carré l'expression de la période : $T^2=4\pi^2\,\dfrac{r^3}{G\,M}$, d'où
\[\frac{T^2}{r^3}=\frac{4\pi^2}{G\,M}.\]
Le rapport ne dépend que de la masse de l'astre attracteur : mesurer la période et le rayon de l'orbite d'un satellite
permet de \textbf{peser} l'astre autour duquel il tourne.
\end{savoirfaire}
\begin{codecote}[0.6]
\begin{pythoncol}
planetes = {"Mercure": (0.387, 0.241),
            "Vénus":   (0.723, 0.615),
            "Terre":   (1.000, 1.000),
            "Mars":    (1.524, 1.881),
            "Jupiter": (5.203, 11.86),
            "Saturne": (9.537, 29.45)}
for nom, (a, T) in planetes.items():
    print(nom, round(T**2 / a**3, 3))
\end{pythoncol}
\begin{consolecol}
Mercure 1.002
Vénus 1.001
Terre 1.0
Mars 1.0
Jupiter 0.999
Saturne 1.0
\end{consolecol}
\cote
\fontsize{9.5}{12.5}\selectfont
\capacite{Capacité numérique}\par\smallskip
Le dictionnaire associe à chaque planète le demi-grand axe $a$ de son orbite, en unités astronomiques, et sa période $T$,
en années. Le rapport $T^2/a^3$ vaut 1 pour toutes les planètes, aux arrondis près : la troisième loi est vérifiée. En
unités SI, il permet de calculer la masse du Soleil : \qty{1,99e30}{kg}.
\end{codecote}
\essSection{4}{Le satellite géostationnaire}
\begin{propriete}[Orbite géostationnaire]
Un satellite \textbf{géostationnaire} reste à la verticale d'un même point de l'équateur : il tourne dans le plan
équatorial, dans le même sens que la Terre, avec une période égale à la durée de rotation de la Terre dans le
référentiel géocentrique, $T=\qty{86164}{s}$ (23~h 56~min). La troisième loi impose alors
\[r=\left(\frac{G\,M_\mathrm{T}\,T^2}{4\pi^2}\right)^{1/3}=\qty{4,215e7}{m},\]
soit une altitude d'environ \num{35800}~km. Les antennes paraboliques qui le visent restent fixes.
\end{propriete}
\begin{attention}
Un satellite en orbite n'est pas « en apesanteur » parce que l'attraction terrestre y serait nulle : à l'altitude de la
station spatiale, elle vaut encore 88~\% de sa valeur au sol. Le satellite et les objets qu'il contient tombent ensemble,
en permanence, autour de la Terre.
\end{attention}
\begin{experience}[Peser Jupiter avec une de ses lunes]
Avec une lunette, ou un logiciel de planétarium, on repère chaque nuit à la même heure la position d'Europe, une lune
de Jupiter, de part et d'autre de la planète. Son orbite étant vue par la tranche, l'élongation $x$ varie de façon
sinusoïdale. La modélisation donne une amplitude de \num{9,39} rayons de Jupiter ($R_\mathrm{J}=\qty{71492}{km}$),
soit $r=\qty{6,71e8}{m}$, et une période $T=\qty{3,551}{j}$. La troisième loi de Kepler conduit alors à
$M_\mathrm{J}=\dfrac{4\pi^2 r^3}{G\,T^2}=\qty{1,90e27}{kg}$.
\end{experience}
\schema{16-europe}
\begin{rediger}
Préciser le référentiel (géocentrique pour un satellite de la Terre) et le système. Écrire la deuxième loi dans le repère
de Frenet, puis projeter : « selon $\ut$ : $\ddt{v}=0$, le mouvement est uniforme ; selon $\un$ : $\frac{v^2}{r}=\frac{GM}{r^2}$ ». Ne
pas confondre l'altitude $h$ et le rayon de l'orbite $r=R+h$.
\end{rediger}
\exercices

\rubrique{Découverte}
\begin{blocLivre}{ex-16-1}
\begin{exercice}[1]{La Terre et la Lune}
La Lune, de masse $M_\mathrm{L}=\qty{7,35e22}{kg}$ et de rayon \qty{1737}{km}, est en orbite autour de la Terre à la distance
$d=\qty{3,84e8}{m}$ (de centre à centre).
\begin{donnees}
$G=\qty{6,67e-11}{N.m^2.kg^{-2}}$ ; $M_\mathrm{T}=\qty{5,97e24}{kg}$.
\end{donnees}
\begin{questions}
\item Calculer la valeur de la force exercée par la Terre sur la Lune. Que vaut la force exercée par la Lune sur la Terre ?
\item En déduire la valeur de l'accélération de la Lune dans le référentiel géocentrique.
\item Calculer l'intensité du champ de gravitation à la surface de la Lune. La comparer à celle de la pesanteur terrestre.
\end{questions}
\end{exercice}
\end{blocLivre}
\begin{blocLivre}{ex-16-2}
\begin{exercice}[1]{Un satellite en orbite basse}
Les satellites d'une constellation de télécommunication orbitent à 550~km d'altitude, sur des orbites supposées
circulaires.
\begin{donnees}
$G\,M_\mathrm{T}=\qty{3,98e14}{m^3.s^{-2}}$ ; $R_\mathrm{T}=\qty{6371}{km}$.
\end{donnees}
\begin{questions}
\item Dans quel référentiel étudie-t-on leur mouvement ? Quelle force s'exerce sur un satellite ?
\item Montrer que le mouvement est circulaire uniforme et établir l'expression de la vitesse.
\item Calculer la vitesse et la période d'un satellite.
\item Un satellite deux fois plus massif aurait-il la même période ? Justifier.
\end{questions}
\end{exercice}
\end{blocLivre}
\begin{blocLivre}{ex-16-3}
\begin{exercice}[1]{L'orbite d'une comète}
Le schéma représente l'orbite elliptique d'une comète autour du Soleil, avec quatre positions occupées à intervalles de
temps égaux.
\schema{16-ex03}
\begin{questions}
\item Énoncer la première loi de Kepler. Où se trouve le Soleil ?
\item En quel point la vitesse de la comète est-elle maximale ? minimale ? Justifier avec la deuxième loi.
\item Comparer les aires des surfaces hachurées. Que peut-on en conclure sur les durées de parcours ?
\item Que deviennent les lois de Kepler pour une orbite circulaire ?
\end{questions}
\end{exercice}
\end{blocLivre}
\begin{blocLivre}{ex-16-4}
\begin{exercice}[1]{Le satellite géostationnaire}
\begin{donnees}
$G\,M_\mathrm{T}=\qty{3,98e14}{m^3.s^{-2}}$ ; $R_\mathrm{T}=\qty{6371}{km}$ ; durée de rotation de la Terre dans le référentiel
géocentrique : \qty{86164}{s}.
\end{donnees}
\begin{questions}
\item Citer les trois conditions que doit remplir un satellite pour être géostationnaire.
\item Pourquoi utilise-t-on \qty{86164}{s} et non \qty{86400}{s} ?
\item À l'aide de la troisième loi de Kepler, calculer le rayon de l'orbite, puis l'altitude.
\item Calculer la vitesse d'un satellite géostationnaire. Pourquoi les antennes de télévision qui le visent sont-elles
fixes ?
\end{questions}
\end{exercice}
\end{blocLivre}
\begin{blocLivre}{ex-16-5}
\begin{exercice}[1]{Peser Jupiter}
Io, l'un des satellites de Jupiter découverts par Galilée en 1610, décrit autour de la planète une orbite quasi circulaire
de rayon \qty{4,22e8}{m}, en 1,769~jour.
\begin{donnees}
$G=\qty{6,67e-11}{N.m^2.kg^{-2}}$.
\end{donnees}
\begin{questions}
\item Dans quel référentiel étudier le mouvement d'Io ?
\item Écrire la troisième loi de Kepler pour une orbite circulaire.
\item En déduire la masse de Jupiter. Combien de fois est-elle plus grande que celle de la Terre
($M_\mathrm{T}=\qty{5,97e24}{kg}$) ?
\end{questions}
\end{exercice}
\end{blocLivre}
\rubrique{Entraînement}
\begin{blocLivre}{ex-16-6}
\begin{exercice}[2]{La comète de Halley}
La comète de Halley passe à \qty{0,586}{ua} du Soleil au périhélie, et s'en éloigne jusqu'à \qty{35,1}{ua} à l'aphélie. Au
périhélie, sa vitesse vaut \qty{54,5}{km.s^{-1}}.
\begin{questions}
\item En utilisant la deuxième loi de Kepler sur une très courte durée au périhélie et à l'aphélie, où la vitesse est
perpendiculaire au rayon, montrer que $r_\mathrm{P}\,v_\mathrm{P}=r_\mathrm{A}\,v_\mathrm{A}$.
\item En déduire la vitesse de la comète à l'aphélie.
\item Calculer le demi-grand axe de l'orbite.
\item À l'aide de la troisième loi, la Terre servant de référence ($a=\qty{1}{ua}$, $T=\qty{1}{an}$), calculer la période de la
comète. Elle est passée au périhélie en 1986 : quand reviendra-t-elle ?
\end{questions}
\end{exercice}
\end{blocLivre}
\begin{blocLivre}{ex-16-7}
\begin{exercice}[2]{Tester la troisième loi en Python}
On reprend le programme de l'essentiel.
\begin{questions}
\item \capacite{Capacité numérique}Quelle structure de données le programme utilise-t-il ? Que contiennent les deux nombres
associés à chaque planète ?
\item Pourquoi le rapport affiché vaut-il 1 pour la Terre ? Interpréter les autres valeurs.
\item Écrire les lignes qui tracent $T^2$ en fonction de $a^3$. Quelle courbe obtient-on ?
\item Convertir le rapport en unités SI ($\qty{1}{ua}=\qty{1,496e11}{m}$ ; $\qty{1}{an}=\qty{3,156e7}{s}$) et en déduire la masse du
Soleil.
\end{questions}
\end{exercice}
\end{blocLivre}
\begin{blocLivre}{ex-16-8}
\begin{exercice}[2]{Peser la Terre grâce à la Lune}
La Lune fait le tour de la Terre en 27,32~jours, dans le référentiel géocentrique, sur une orbite que l'on supposera
circulaire de rayon \qty{3,844e5}{km}.
\begin{donnees}
$G=\qty{6,67e-11}{N.m^2.kg^{-2}}$.
\end{donnees}
\begin{questions}
\item Établir l'expression de la vitesse de la Lune en fonction de $G$, $M_\mathrm{T}$ et $r$.
\item En déduire l'expression de la période, puis la troisième loi de Kepler.
\item Calculer la masse de la Terre.
\item Comparer à la valeur $\qty{5,97e24}{kg}$ et proposer une raison à l'écart.
\end{questions}
\end{exercice}
\end{blocLivre}
\begin{blocLivre}{ex-16-9}
\begin{exercice}[2]{Les satellites GPS}
Les satellites du système GPS évoluent sur des orbites circulaires à \num{20200}~km d'altitude.
\begin{donnees}
$G\,M_\mathrm{T}=\qty{3,98e14}{m^3.s^{-2}}$ ; $R_\mathrm{T}=\qty{6371}{km}$.
\end{donnees}
\begin{questions}
\item Calculer la période de révolution d'un satellite GPS. La comparer à la durée de rotation de la Terre.
\item Calculer sa vitesse.
\item Un satellite GPS est-il géostationnaire ? Combien de fois repasse-t-il au-dessus du même point du globe en 24~h ?
\end{questions}
\end{exercice}
\end{blocLivre}
\begin{blocLivre}{ex-16-10}
\begin{exercice}[2]{Observer la Terre}
Les satellites d'observation de la Terre de la série Sentinel-2 évoluent à 786~km d'altitude, sur une orbite circulaire
passant presque au-dessus des pôles.
\begin{donnees}
$G\,M_\mathrm{T}=\qty{3,98e14}{m^3.s^{-2}}$ ; $R_\mathrm{T}=\qty{6371}{km}$.
\end{donnees}
\begin{questions}
\item Calculer la période de révolution et le nombre de tours effectués par jour.
\item Pourquoi une orbite polaire permet-elle d'observer toute la surface de la Terre au fil des jours ?
\item Un tel satellite pourrait-il être géostationnaire ? Justifier.
\end{questions}
\end{exercice}
\end{blocLivre}
\begin{blocLivre}{ex-16-11}
\begin{exercice}[2]{Lancer depuis Kourou}
Pour être mis en orbite circulaire très basse, un satellite doit atteindre la vitesse $v_1=\sqrt{G\,M_\mathrm{T}/R_\mathrm{T}}$. La
rotation de la Terre donne au lanceur, avant même son décollage, une vitesse vers l'est qui dépend de la latitude
$\lambda$ du site : $v_\mathrm{rot}=\frac{2\pi R_\mathrm{T}}{T}\cos\lambda$.
\begin{donnees}
$G\,M_\mathrm{T}=\qty{3,98e14}{m^3.s^{-2}}$ ; $R_\mathrm{T}=\qty{6371}{km}$ ; $T=\qty{86164}{s}$ ; Kourou (Guyane) : $\lambda=5{,}2^\circ$ ;
Baïkonour (Kazakhstan) : $\lambda=46^\circ$.
\end{donnees}
\begin{questions}
\item Calculer $v_1$.
\item Calculer la vitesse apportée par la rotation de la Terre à Kourou et à Baïkonour.
\item Exprimer chacune en pourcentage de $v_1$. Pourquoi le site de Kourou est-il avantageux ?
\item Pourquoi lance-t-on les satellites vers l'est ?
\end{questions}
\end{exercice}
\end{blocLivre}
\rubrique{Objectif bac}
\begin{blocLivre}{ex-16-12}
\begin{exercice}[3]{Le télescope spatial Hubble}
\begin{docu}{1}{Un observatoire en orbite}
Lancé en 1990, le télescope spatial Hubble orbite autour de la Terre à environ 540~km d'altitude, sur une orbite
quasi circulaire. Au-dessus de l'atmosphère, il échappe à la turbulence qui brouille les images des télescopes au sol.
\end{docu}
\begin{docu}{2}{Une orbite qui se dégrade}
L'atmosphère très ténue qui subsiste à cette altitude freine légèrement le télescope : son orbite s'abaisse lentement,
et des missions de navettes l'ont plusieurs fois remonté.
\end{docu}
\begin{donnees}
$G\,M_\mathrm{T}=\qty{3,98e14}{m^3.s^{-2}}$ ; $R_\mathrm{T}=\qty{6371}{km}$.
\end{donnees}
\begin{questions}
\item Justifier que le mouvement de Hubble est circulaire uniforme, puis calculer sa vitesse et sa période.
\item Représenter, sur un schéma, les vecteurs vitesse et accélération de Hubble.
\item \initiative Si l'orbite s'abaisse à 500~km, comment évoluent la vitesse et la période ? Calculer leurs nouvelles valeurs.
\item Commenter ce résultat paradoxal : le frottement de l'air finit par augmenter la vitesse du télescope.
\end{questions}
\end{exercice}
\end{blocLivre}
\begin{blocLivre}{ex-16-13}
\begin{exercice}[3]{Les deux lunes de Mars}
\begin{docu}{1}{Phobos et Deimos}
Mars possède deux petites lunes. Phobos orbite à \qty{9376}{km} du centre de Mars, en 7,66~h ; Deimos, à \qty{23463}{km}, en
30,3~h. Leurs orbites sont quasi circulaires et situées dans le plan de l'équateur martien.
\end{docu}
\begin{docu}{2}{Une lune qui se lève à l'ouest}
Mars tourne sur elle-même en 24,6~h. Vu depuis la surface de Mars, Phobos se lève à l'ouest et se couche à l'est, deux
fois par jour.
\end{docu}
\begin{donnees}
$G=\qty{6,67e-11}{N.m^2.kg^{-2}}$.
\end{donnees}
\begin{questions}
\item Établir la troisième loi de Kepler pour une orbite circulaire.
\item Calculer la masse de Mars à partir des données de Phobos, puis de Deimos. Commenter.
\item \initiative Calculer le rayon de l'orbite « aréostationnaire », sur laquelle un satellite resterait à la verticale
d'un point de l'équateur martien.
\item Expliquer pourquoi Phobos semble se lever à l'ouest, à la différence de la Lune vue depuis la Terre.
\end{questions}
\end{exercice}
\end{blocLivre}
\begin{blocLivre}{ex-16-14}
\begin{exercice}[3]{Les sept planètes de TRAPPIST-1}
\begin{docu}{1}{Un système planétaire voisin}
En 2017, l'annonce de sept planètes de la taille de la Terre autour de l'étoile TRAPPIST-1, à 40 années-lumière, a fait
sensation. Leurs périodes ont été mesurées par la méthode des transits, leurs distances déduites des observations.
\par\smallskip
{\small\renewcommand{\arraystretch}{1.2}
\begin{tabular}{@{}lccc@{}}\toprule
Planète & b & e & h\\\midrule
demi-grand axe $a$ (ua) & \num{0.01154} & \num{0.02925} & \num{0.06189}\\
période $T$ (jours) & \num{1.511} & \num{6.100} & \num{18.77}\\\bottomrule
\end{tabular}}
\end{docu}
\begin{donnees}
Pour le système solaire, la troisième loi s'écrit $\dfrac{T^2}{a^3}=1$, avec $T$ en années et $a$ en ua. On admet que, pour une étoile
de masse $M$, $\dfrac{T^2}{a^3}=\dfrac{M_\odot}{M}$ avec les mêmes unités, $M_\odot$ étant la masse du Soleil. 1~an $=$ 365,25~jours.
\end{donnees}
\begin{questions}
\item Énoncer la troisième loi de Kepler. Pourquoi le rapport $T^2/a^3$ dépend-il de l'étoile ?
\item Calculer le rapport $T^2/a^3$, en an$^2$/ua$^3$, pour les trois planètes. Commenter.
\item \initiative En déduire la masse de l'étoile TRAPPIST-1 en masses solaires.
\item La planète e reçoit de son étoile une puissance par unité de surface proche de celle que la Terre reçoit du Soleil.
Comment est-ce possible alors qu'elle est 34 fois plus proche de son étoile ?
\end{questions}
\end{exercice}
\end{blocLivre}
\begin{blocLivre}{ex-16-15}
\begin{exercice}[3]{L'orbite de transfert}
\begin{docu}{1}{Rejoindre l'orbite géostationnaire}
Un satellite placé sur une orbite circulaire basse, de rayon $r_1=\qty{6671}{km}$, doit rejoindre l'orbite géostationnaire
de rayon $r_2=\qty{42150}{km}$. On l'accélère brièvement : il décrit alors une demi-ellipse, dont le périgée est sur l'orbite
basse et l'apogée sur l'orbite géostationnaire, où une seconde poussée l'y maintient.
\end{docu}
\schema{16-ex15}
\begin{donnees}
$G\,M_\mathrm{T}=\qty{3,98e14}{m^3.s^{-2}}$.
\end{donnees}
\begin{questions}
\item Calculer les vitesses du satellite sur l'orbite basse et sur l'orbite géostationnaire.
\item Quelle est la valeur du demi-grand axe de l'ellipse de transfert ?
\item \initiative En admettant que la troisième loi de Kepler s'applique à l'ellipse avec son demi-grand axe, calculer la
durée du transfert.
\item Pourquoi la vitesse du satellite diminue-t-elle pendant le transfert ? Justifier avec la deuxième loi de Kepler.
\end{questions}
\end{exercice}
\end{blocLivre}
